% texte.tex — f176r, vv. 1128–1168 (conversion initiale des travaux antérieurs)
\vers[n=1127]{}{}{τὸ ξιπομάκαιραν ἀποκεκόψι τουτοΐ.}{avec ce sabre.}

\vers[n=1128]{ΕΥ.}{Euripide}{αἰαῖ· τί δράσω; πρὸς τίνας στρεφθῶ λόγους;}{Aïe, aïe! que faire ? à quelles raisons recourir ?}
\vers[n=1129]{}{}{ἀλλ᾽ οὐκ ἂν ἐνδέξαιτο βάρβαρος φύσις.}{Mais cette nature barbare ne les admettrait pas.}
\vers[n=1130]{}{}{σκαιοῖσι γάρ τοι καινὰ προσφέρων σοφὰ}{« Proposer à des sots des doctrines nouvelles,}
\vers[n=1131]{}{}{μάτην ἀναλίσκοις ἄν, ἀλλ᾽ ἄλλην τινὰ}{c'est perdre sa peine. » Eh bien, dressons quelqu'autre}
\vers[n=1132]{}{}{τούτῳ πρέπουσαν μηχανὴν προσοιστέον.}{machine, qui soit de sa portée. (Il s'en va.)}
\vers[n=1133]{ΤΟ.}{Archer}{μιαρὸς ἀλώπηξ, οἶον ἐπιτήκιζί μοι.}{La canaille de renard, comme il pensait m'embabouiner !}
\vers[n=1134]{ΜΝ.}{Mnésiloque}{μέμνησο Περσεῦ μ᾽ ὡς καταλείπεις ἀθλίαν.}{Souviens-toi, Persée, que tu me laisses en vilaine posture.}
\vers[n=1135]{ΤΟ.}{Archer}{ἔτι γὰρ σὺ τὴ μάστιγαν ἐπιτυμεῖς λαβεῖν;}{As-tu encore envie d'attraper le fouet? (Il se recouche et s'endort.)}
\vers[n=1136, indent=2em, findent=1.2em, metre={lyr.}]{ΧΟ.}{Chœur}{Παλλάδα τὴν φιλόχορον ἐμοὶ}{C'est ma coutume d'appeler parmi nous Pallas,}
\vers[n=1137, indent=2em, findent=1.2em]{}{}{δεῦρο καλεῖν νόμος εἰς χορόν,}{amie du chœur, la jeune fille,}
\vers[n=1138, indent=2em, findent=1.2em]{}{}{παρθένον ἄζυγα κούρην,}{la vierge libre du joug,}
\vers[n=1140, indent=2em, findent=1.2em]{}{}{ἣ πόλιν ἡμετέραν ἔχει}{gardienne de notre acropole,}
\vers[n=1141, indent=2em, findent=1.2em]{}{}{καὶ κράτος φανερὸν μόνη}{seule en possession d'un visible pouvoir,}
\vers[n=1142, indent=2em, findent=1.2em]{}{}{κλῃδοῦχός τε καλεῖται.}{et qui a nom Clédouque.}
\vers[n=1143, indent=2em, findent=1.2em]{}{}{φάνηθ᾽ ὦ τυράννους στυγοῦσ᾽ ὥσπερ εἰκός.}{Apparais, toi qui détestes les tyrans, comme faire se doit.}
\vers[n=1145, indent=2em, findent=1.2em]{}{}{δῆμός τοί σε καλεῖ γυναι-}{Le peuple des femmes t'appelle.}
\vers[n=1146, indent=2em, findent=1.2em]{}{}{\alinea κῶν· ἔχουσα δέ μοι μόλοις}{Viens à moi avec}
\vers[n=1147, indent=2em, findent=1.2em]{}{}{εἰρήνην φιλέορτον.}{la Paix amie des fêtes.}
\vers[n=1148, indent=2em, findent=1.2em]{}{}{ἥκετέ τ᾽ εὔφρονες ἵλαοι,}{Venez, augustes déesses, bienveillantes et propices,}
\vers[n=1149, indent=2em, findent=1.2em]{}{}{πότνιαι, ἄλσος ἐς ὑμέτερον,}{dans votre enceinte sacrée,}
\vers[n=1150, indent=2em, findent=1.2em]{}{}{οὗ δὴ ἀνδράσιν οὐ θεμιτεύετον εἰσορᾶν}{où vous ne permettez pas aux hommes de contempler}
\vers[n=1151, indent=2em, findent=1.2em]{}{}{ὄργια σεμνὰ θεαῖν ἵνα λαμπάσι.}{vos saintes orgies à la lueur des torches.}
\vers[n=1152, indent=2em, findent=1.2em]{}{}{φαίνετον ἄμβροτον ὄψιν.}{Découvrez vos faces immortelles.}
\vers[n=1153, indent=2em, findent=1.2em]{}{}{μόλετον, ἔλθετον, ἀντόμεθ᾽ ὦ}{Venez, approchez, nous vous en conjurons,}
\vers[n=1154, indent=2em, findent=1.2em]{}{}{Θεσμοφόρω πολυποτνία,}{Thesmophores très vénérées.}
\vers[n=1155, indent=2em, findent=1.2em]{}{}{εἰ καὶ πρότερόν ποτ᾽ ἐπηκόω ἤλθετον}{Si jamais auparavant vous répondîtes à notre appel,}
\vers[n=1156, indent=2em, findent=1.2em]{}{}{νῦν ἀφίκεσθον †ἱκετεύομεν ἐνθάδ᾽ ἡμῖν†.}{venez aujourd'hui, sur nos instantes prières, ici au milieu de nous.}
\vers[n=1160]{ΕΥ.}{Euripide}{γυναῖκες εἰ βούλεσθε τὸν λοιπόν χρόνον}{(Déguisé en vieille procureuse et portant une pectis ou harpe lydienne. Il est accompagné d'une danseuse et d'un jeune flûtiste). Femmes, si vous voulez pour le reste du temps}
\vers[n=1161]{}{}{σπονδὰς ποιήσασθαι πρὸς ἐμέ, νυνὶ πάρα,}{faire la paix avec moi, c'est le moment,}
\vers[n=1162]{}{}{ἐφ᾽ ᾧτ᾽ ἀκοῦσαι μηδὲν ὑπ᾽ ἐμοῦ μηδαμὰ}{à charge que jamais plus à l'avenir je ne dirai de vous}
\vers[n=1163]{}{}{κακὸν τὸ λοιπόν. ταῦτ᾽ ἐπικηρυκεύομαι.}{le moindre mal. Telles sont mes propositions.}
\vers[n=1164]{ΚΟ.}{Coryphée}{χρείᾳ δὲ ποίᾳ τόνδ᾽ ἐπεσφέρεις λόγον;}{Qu'est-ce qui te décide à faire ces ouvertures ?}
\vers[n=1165]{ΕΥ.}{Euripide}{ὅδ᾽ ἐστὶν οὑν τῇ σανίδι κηδεστὴς ἐμός.}{Cet homme attaché à l'ais est mon beau-père.}
\vers[n=1166]{}{}{ἢν οὖν κομίσωμαι τοῦτον, οὐδὲν μή ποτε}{Que je le recouvre, et jamais}
\vers[n=1167]{}{}{κακῶς ἀκούσητ᾽· ἢν δὲ μὴ πίθησθέ μοι,}{je ne médirai de vous. Si vous refusez,}
\vers[n=1168]{}{}{ἃ νῦν ὑποικουρεῖτε τοῖσιν ἀνδράσιν}{toutes vos pratiques à huis-clos, je les révélerai}
